% Größen und Einheiten – bank/einheiten/ (zusammenbau v0.4, Heft)
\einheitenkopf[t7]{Größen und Einheiten}
\setcounter{aufgabe}{20}

\begin{aufgabe}{Länge, Masse und Geld umrechnen – Prüfungsaufgaben}
\begin{teile}
\teil Setze $<$, $>$ oder $=$ ein: $0{,}04$ m \leerfeld $40$ cm \hfill (P10 2019 OS)
\teil Setze $<$, $=$ oder $>$ ein: $2{,}7$ m \leerfeld $27$ cm \hfill (P10 2026 FOR)
\teil Setze $<$, $>$ oder $=$ ein: $0{,}07$ kg \leerfeld $70$ g \hfill (P10 2019 OS)
\end{teile}
\end{aufgabe}

\begin{aufgabe}{Zeit umrechnen und Zeitspannen berechnen – Prüfungsaufgaben}
\begin{teile}
\teil Eine Radtour dauert $5{,}5$ h. Wie viele Minuten sind das? \hfill (P10 2018 OS) \\ \leerfeld[min]
\teil Eine Schicht dauert $6{,}5$ Stunden. Gib die Dauer in Minuten an. \hfill (P10 2025 OS) \\ \leerfeld[min]
\teil Wie viele Minuten sind $2{,}25$ Stunden? Kreuze an. \hfill (P10 2024 OS)\\ \kreuz{$125$ min}\\ \kreuz{$135$ min}\\ \kreuz{$145$ min}\\ \kreuz{$155$ min}
\teil Ein Zug fährt um 9:47 Uhr ab und ist $3$ h $25$ min unterwegs. Wann kommt er an? \hfill (P10 2021 OS) \\ \leerfeld Uhr
\teil Ein Fernbus fährt um 16:52 Uhr ab, die Fahrt dauert $1$ h $43$ min. Wann kommt er an? \hfill (P10 2021 OS) \\ \leerfeld Uhr
\end{teile}
\end{aufgabe}

\begin{aufgabe}{Flächen- und Volumeneinheiten umrechnen – Prüfungsaufgaben}
\begin{teile}
\teil Eine Flasche enthält $1{,}25$ l Saft. Wie viele Gläser zu $200$ ml kann man damit ganz füllen? \hfill (P10 2022 OS) \\ \leerfeld Gläser
\teil Ein Kanister enthält $4{,}2$ l Wasser. Wie viele Trinkflaschen zu $750$ ml kann man damit ganz füllen? \hfill (P10 2022 OS) \\ \leerfeld Flaschen
\teil In einem Topf sind $2{,}7$ l Suppe. Wie viele volle Kellen zu $250$ ml kann man ausgeben? \hfill (P10 2022 OS) \\ \leerfeld Kellen
\end{teile}
\end{aufgabe}
